\begin{textsample}{POS Dim 6 – global\_north\_2024\_07 – Score 6.00 – global\_north\_2024\_07/110372368.txt}  \label{ex:f6_pos_235}
It was never touted as a complete solution to the problems of getting humanitarian aid into Gaza.
Now, the United States has announced its temporary aid pier will wind down operations after being plagued with challenges and bad weather.
The United Nations has long said \textbf{maritime} deliveries were no substitute for land access, saying land routes need to remain the focus of aid operations.
Here ' s how the pier was supposed to work -- and why it faced multiple setbacks.
Logistics and security of pier were complicated In March, US President Joe Biden used his State of the Union address to announce the military would begin constructing the pier that would receive large \textbf{shipments} of food, water, medicine and temporary shelter.
Despite there being"no US boots on the ground,"Mr Biden said the pier would"enable a massive increase in the amount of humanitarian assistance getting into Gaza every day".
Loading...
The Pentagon said the system was expected to cost at least \$US320 million ( \$490 million ).
The logistics @ @ @ @ @ @ @ @ @ @ complicated : Since Israel had concerns that \textbf{shipments} would contain items Hamas could use against its troops, all humanitarian aid had to be sent first to Cyprus for security checks at Larnaca port.
Once pallets were inspected, they were loaded onto \textbf{ships} -- mainly commercial \textbf{vessels} -- and taken about 321 kilometres to the large floating pier.
Palettes were then transferred onto trucks that were then loaded onto two types of smaller army \textbf{boats}.
The trucks were then shuttled from the pier to a floating causeway located several miles away that was anchored into the beach by Israel Defense Forces.
Aid groups then collected the supplies for distribution on shore at a port facility built by the Israelis just south-west of Gaza City.
There were three zones to the port : 1.
One controlled by the Israelis where aid from the pier was dropped off, 2.
One where the aid will be transferred, 3.
A third where Palestinian drivers contracted by the UN will wait to pick up the @ @ @ @ @ @ @ @ @ @ Israel must do its part,"Mr Biden said."Israel must allow more aid into Gaza and ensure humanitarian works aren't caught in the crossfire."Rough \textbf{seas} broke pier almost instantly It took nearly two months for the US to build the pier after Mr Biden announced the plan in his national address.
But in late May, just two weeks after its completion, it broke apart due to rough \textbf{seas} brought on by stormy weather.
Satellite imagery captured a view of the remaining section of the pier after it broke in rough \textbf{seas}. ( Maxar Technologies via AP ) It took at least two days to be pulled from the beach to be taken for repairs, which lasted over a week."From when it was operational, it was working,"Pentagon spokeswoman Sabrina Singh said at the time."We just had... an unfortunate confluence of... storms that made it inoperable for a bit."During that incident, three service @ @ @ @ @ @ @ @ @ @ a critical condition.
The pause came after the Israeli military used an area near the pier to fly out hostages after their rescue in a raid that killed more than 270 Palestinians, prompting a UN security review over concerns for aid workers ' safety.
Israel said it allowed hundreds of truckloads of aid through the crossing, but that the UN had failed to pick it up.
The UN said it is too dangerous for trucks to move through the area, despite Israeli pledges to carve out a safe corridor.
The pier faced criticism from the start Aid groups have regularly criticised the plan to deliver aid to Gaza by \textbf{sea} as ineffective.
The head of the World Health Organization ' s Eastern Mediterranean region said the pier couldn't supply Palestinians with anywhere near the level of aid they needed."The pier has supported a little bit, but it ' s not to the scale that is needed by any stretch of the imagination,"Hanan Balkhy told The Associated Press. @ @ @ @ @ @ @ @ @ @ the end of June, Ms Singh said land routes were the most effective delivery method."We continue to urge for those land routes to be reopened,"she said.
White House National Security Adviser Jake Sullivan said that the pier had helped bring urgently needed food and humanitarian aid to Gaza, but there were now additional supplies entering the Palestinian enclave via land routes."The real issue right now is not about getting aid into Gaza.
It ' s about getting aid around Gaza effectively,"he said.
UN spokesperson Stephane Dujarric said the pier was a welcomed additional resource while it worked."We will keep pushing for what we actually need, which is large-scale road transfer of aid into Gaza,"he said.

% matched lemmas: boat, maritime, sea, ship, shipment, vessel
\end{textsample}
